\documentclass[11pt,a4paper]{article}
\usepackage[utf8]{inputenc}
\usepackage{amsmath,amssymb,amsthm}
\usepackage{algorithm}
\usepackage{algpseudocode}
\usepackage{booktabs}
\usepackage{hyperref}
\usepackage{geometry}
\geometry{margin=1in}

\title{Beam Search with Length Normalization: Heuristic Tree Pruning, Length Bias Penalization, and Decoding Complexity}
\author{Automorphic Intelligence Architecture Specifications}
\date{\today}

\newtheorem{theorem}{Theorem}
\newtheorem{lemma}[theorem]{Lemma}
\newtheorem{proposition}[theorem]{Proposition}
\theoremstyle{definition}
\newtheorem{definition}{Definition}
\newtheorem{remark}{Remark}

\begin{document}

\maketitle

\begin{abstract}
Exact maximum a posteriori (MAP) decoding of autoregressive sequences is NP-hard due to exponential vocabulary branching $\mathcal{O}(|\mathcal{V}|^T)$. Greedy decoding provides an $\mathcal{O}(T)$ approximation but readily collapses into sub-optimal local extrema. Beam search strikes an optimal Pareto efficiency between computational tractability and search completeness by maintaining a bounded frontier of the top-$B$ partial sequence hypotheses. Furthermore, because cumulative log-probabilities decrease monotonically with length, unnormalized beam search systematically favors short sequences. Length normalization penalties resolve this length bias. This document formalizes the pruned beam search tree, analyzes length penalty bounds, details algorithmic pseudocode, and provides a verified step-by-step numerical calculation.
\end{abstract}

\section{Notation and Mathematical Preliminaries}

Let $\mathcal{V}$ denote the discrete token vocabulary, $B \in \mathbb{N}_{\ge 1}$ denote the beam width, and $\mathcal{H}_t = \{(\mathbf{y}_t^{(1)}, s_t^{(1)}), \dots, (\mathbf{y}_t^{(B)}, s_t^{(B)})\}$ denote the active hypothesis set at time step $t$.

\begin{table}[h]
\centering
\caption{Mathematical Notation for Beam Search Decoding}
\begin{tabular}{lll}
\toprule
\textbf{Symbol} & \textbf{Domain} & \textbf{Semantic Description} \\
\midrule
$B$ & $\mathbb{N}_{\ge 1}$ & Beam search width capacity (number of concurrent active hypotheses) \\
$\mathbf{y} = (y_1, \dots, y_L)$ & $\mathcal{V}^L$ & Candidate generated token sequence \\
$s(\mathbf{y})$ & $\mathbb{R}_{\le 0}$ & Cumulative unnormalized log-probability: $\sum_{t=1}^L \log P(y_t \mid y_{<t})$ \\
$\text{LP}(L)$ & $\mathbb{R}_{> 0}$ & Length penalty normalization factor \\
$\alpha$ & $[0, 1.5]$ & Length normalization exponent strength ($\alpha = 0.6$ standard GNMT) \\
$\mathcal{S}_{\text{norm}}(\mathbf{y})$ & $\mathbb{R}$ & Length-normalized hypothesis ranking score \\
$\mathcal{F}$ & Set & Finished set of completed hypotheses terminated by $\langle \text{EOS} \rangle$ \\
\bottomrule
\end{tabular}
\end{table}

\section{Problem Definition and Core Concepts}

\subsection{3-Level Explanation}
\begin{enumerate}
    \item \textbf{Intuitive (High School level):} A greedy decoder picks only the single best word right now, which might lead to a dead end. Full search explores all quadrillion combinations, which takes forever. Beam search explores the top $B$ most promising parallel storylines at every step, dividing the final score by story length so longer sentences aren't unfairly penalized.
    \item \textbf{Engineering Level:} Beam search expands $B$ active sequences by $|\mathcal{V}|$ tokens each step, computing $B \times |\mathcal{V}|$ candidate log-probs, sorting, and retaining the top $B$. Hypotheses ending with $\langle \text{EOS} \rangle$ are moved to a finished buffer and ranked using $\text{Score} = \frac{\log P(\mathbf{y})}{\text{LP}(|\mathbf{y}|)}$.
    \item \textbf{Mathematical Level:} Given sequence model $P_\theta$, beam search approximates $\mathbf{y}^* = \arg\max_\mathbf{y} \mathcal{S}_{\text{norm}}(\mathbf{y})$:
    \begin{equation}
    \mathcal{S}_{\text{norm}}(\mathbf{y}) = \frac{\sum_{t=1}^{|\mathbf{y}|} \log P_\theta(y_t \mid y_{<t}, \mathbf{x})}{\text{LP}(|\mathbf{y}|)}, \quad \text{LP}(L) = \left( \frac{5 + L}{6} \right)^\alpha
    \end{equation}
\end{enumerate}

\subsection{5-Step Algorithmic Framework}
\begin{enumerate}
    \item \textbf{Beam Initialization:} Initialize frontier $\mathcal{H}_0 = \{([\langle \text{BOS} \rangle], 0.0)\}$ and completed set $\mathcal{F} = \emptyset$.
    \item \textbf{Hypothesis Expansion:} For each active beam, evaluate model log-probabilities across all $|\mathcal{V}|$ tokens.
    \item \textbf{Frontier Pruning:} Select top-$B$ candidates with highest cumulative raw log-probabilities.
    \item \textbf{EOS Separation:} Partition newly pruned candidates: routes ending in $\langle \text{EOS} \rangle$ enter $\mathcal{F}$; others form next frontier $\mathcal{H}_{t+1}$.
    \item \textbf{Length-Normalized Re-ranking:} Divide cumulative scores of all candidates in $\mathcal{F}$ by $\text{LP}(L)$ and return top ranked sequence.
\end{enumerate}

\section{Algorithm Specification}

\begin{algorithm}[H]
\caption{Beam Search with GNMT Length Normalization}
\label{alg:beam_search}
\begin{algorithmic}[1]
\Require Beam width $B$, max steps $T_{\max}$, penalty exponent $\alpha$, model transition $P(\cdot \mid \mathbf{y})$.
\Ensure Optimal sequence $\mathbf{y}^* \in \mathcal{V}^*$.
\State $\mathcal{H} \gets \{([\langle \text{BOS} \rangle], 0.0)\}, \quad \mathcal{F} \gets \emptyset$
\For{$t = 1$ \textbf{to} $T_{\max}$}
    \State $\mathcal{C} \gets \emptyset$
    \For{$(\mathbf{y}, s) \in \mathcal{H}$}
        \If{$\mathbf{y}[-1] = \langle \text{EOS} \rangle$}
            \State $\mathcal{F} \gets \mathcal{F} \cup \{(\mathbf{y}, s)\}$
            \State \textbf{continue}
        \EndIf
        \State $\mathbf{p} \gets \log P(\cdot \mid \mathbf{y})$
        \For{$v \in \mathcal{V}$}
            \State $\mathcal{C} \gets \mathcal{C} \cup \{(\mathbf{y} \circ [v], s + \mathbf{p}[v])\}$
        \EndFor
    \EndFor
    \State $\mathcal{H} \gets \text{Top-B}_{\text{score}}(\mathcal{C})$
    \If{$\mathcal{H} = \emptyset$} \textbf{break} \EndIf
\EndFor
\State $\mathcal{F} \gets \mathcal{F} \cup \mathcal{H}$
\State $\mathbf{y}^* \gets \arg\max_{(\mathbf{y}, s) \in \mathcal{F}} \frac{s}{((5 + |\mathbf{y}|)/6)^\alpha}$
\State \Return $\mathbf{y}^*$
\end{algorithmic}
\end{algorithm}

\section{Theoretical Guarantees and Search Bounds}

\begin{theorem}[Monotonic Score Non-Increase and Sub-Optimality Bound]
\label{thm:beam_search_bound}
Since $\log P(y_t \mid y_{<t}) \le 0$, the raw cumulative score $s(\mathbf{y})$ is monotonically non-increasing in length: $s(\mathbf{y} \circ [v]) \le s(\mathbf{y})$. The beam search return satisfies:
\begin{equation}
s(\mathbf{y}_{\text{greedy}}) \le s(\mathbf{y}_{\text{beam}}) \le s(\mathbf{y}_{\text{exact MAP}}^*)
\end{equation}
\end{theorem}

\begin{proof}
Probabilities are bounded $P \in (0, 1] \implies \log P \in (-\infty, 0]$. Adding a negative value strictly decreases the sum. Greedy search is the degenerate special case of beam search with $B=1$, while exact MAP corresponds to $B = |\mathcal{V}|^{T_{\max}}$. Since the candidate search space monotonically expands with $B$, the maximum score over the frontier increases monotonically with $B$.
\end{proof}

\section{Complexity Analysis}

\begin{itemize}
    \item \textbf{Time Complexity:} $\mathcal{O}(T \cdot (B \cdot \text{Cost}_{\text{model}} + B |\mathcal{V}| \log B))$ total operations.
    \item \textbf{Space Complexity:} $\mathcal{O}(B \cdot T)$ working memory to store active token sequences.
\end{itemize}

\section{Exactness and Error Analysis}

Without length normalization ($\alpha = 0$), raw score $s(\mathbf{y}) = \sum_{t=1}^L \log P_t \le L \log P_{\max}$. For $P_{\max} < 1$, the score decays linearly with length $\mathcal{O}(L)$, introducing an exponential bias towards prematurely short outputs. Setting $\alpha \in [0.6, 0.8]$ cancels this length penalty.

\section{Worked Numerical Example}

Let $B=2, \alpha = 1.0, \mathcal{V} = \{A, B, \text{EOS}\}$.
Step 1: Expansions from BOS:
\begin{itemize}
    \item $[\text{BOS}, A]$: score $\log(0.6) = -0.5108$.
    \item $[\text{BOS}, B]$: score $\log(0.3) = -1.2040$.
    \item $[\text{BOS}, \text{EOS}]$: score $\log(0.1) = -2.3026$.
    \item Beam selects Top-2: $\{([\text{BOS}, A], -0.5108), ([\text{BOS}, B], -1.2040)\}$.
\end{itemize}
Step 2: Expand from $[\text{BOS}, A]$:
\begin{itemize}
    \item $[\text{BOS}, A, \text{EOS}]$: score $-0.5108 + \log(0.8) = -0.5108 - 0.2231 = -0.7339$.
    \item $[\text{BOS}, A, A]$: score $-0.5108 + \log(0.2) = -0.5108 - 1.6094 = -2.1202$.
\end{itemize}
Completed candidate $[\text{BOS}, A, \text{EOS}]$ (length $L=2$).
With $\alpha = 0.6$: $\text{LP}(2) = ((5+2)/6)^{0.6} = (7/6)^{0.6} = (1.1667)^{0.6} \approx 1.097$.
Normalized score $= \frac{-0.7339}{1.097} \approx -0.6690$.

\section{Numerical Considerations and Edge Cases}

\begin{itemize}
    \item \textbf{Diverse Beam Search:} Standard beam search frequently produces near-identical beams (differing by a single punctuation mark). Grouping beams into diverse subsets with inter-group Hamming penalties enhances output diversity.
\end{itemize}

\section{References}
\begin{enumerate}
    \item Wu, Y., Schuster, M., Chen, Z., Le, Q. V., Norouzi, M., Macherey, W., ... & Dean, J. (2016). Google's neural machine translation system: Bridging the gap between human and machine translation. \emph{arXiv preprint arXiv:1609.08144}.
    \item Vijayakumar, A. K., Cogswell, M., Selvaraju, R. R., Sun, Q., Lee, S., Crandall, D., & Batra, D. (2018). Diverse beam search for out-of-distribution sequence generation. \emph{IJCV}, 126(8), 869--884.
\end{enumerate}

\end{document}
